%%%%%%%%%%%%%%%%%%%%%%%%%%%%%%%%%%%%%%%%%%%%%%%%%%%%%%%%%%%%%%%%%%%%%%%%%%%%%%%%%%%%%%%%%%%%%%%%%%%%%%
%
%   Filename    : abstract.tex 
%
%   Description : This file will contain your abstract.
%                 
%%%%%%%%%%%%%%%%%%%%%%%%%%%%%%%%%%%%%%%%%%%%%%%%%%%%%%%%%%%%%%%%%%%%%%%%%%%%%%%%%%%%%%%%%%%%%%%%%%%%%%

\begin{abstract}
The Taal Volcano in the Philippines remains one of the country's most active volcanoes due to its recurring eruptive activity, highlighting the continued need for effective tools that support hazard communication, disaster preparedness, and scientific visualization.
Interactive three-dimensional computer graphics simulations have demonstrated their potential to provide intuitive visual representations of volcanic eruption plumes while allowing users to explore eruption scenarios in an immersive virtual environment.
Existing graphics-based volcanic plume simulators have successfully demonstrated real-time visualization capabilities; however, many rely on legacy graphics rendering technologies and software architectures that limit extensibility, maintainability, and the efficient utilization of modern graphics hardware. 
This study proposes a modern 3D graphics rendering framework for interactive Taal Volcano plume simulation by redesigning the rendering architecture of an existing plume visualization system while preserving its established particle-based simulation model and core visualization techniques.
The proposed framework aims to modernize the rendering pipeline, graphics resource management, shader implementation, and software architecture to provide a more scalable, maintainable, and extensible platform for future graphics-based scientific visualization applications.
A design and development research methodology will be employed to reimplement the simulator using a contemporary graphics rendering framework and evaluate its rendering performance, graphics resource utilization, scalability, and visual fidelity through quantitative performance metrics and qualitative visual assessment.


%%\begin{itemize}
  % \item \textbf{Motivation}. One to two sentence(s) describing the research domain or problem area.
   %\item \textbf{Problem}. One to two sentence(s) describing the research challenge or opportunity to be addressed in your study
   %\item \textbf{Contribution}. One to two sentence(s) describing your intended contribution or how you plan to address the research question 
   %\item \textbf{Methodology}. One to two sentence(s) summarizing the approach or the manner in which you will carry out your study
   %\item \textbf{Significant Results or Findings}. One to two sentence(s) describing your findings and result (This can be skipped during proposal stage) 
%%\end{itemize}


%
%  Do not put citations or quotes in the abract.
%


\begin{flushleft}
\begin{tabular}{lp{4.25in}}
\hspace{-0.5em}\textbf{Keywords:}\hspace{0.25em} & Computer graphics, 3D graphics, Graphics rendering, Volcanic plume simulation, Taal Volcano, Particle systems
\end{tabular}
\end{flushleft}
\end{abstract}
